\documentclass[11pt]{article}
\usepackage[margin=1in]{geometry}
\usepackage{amsmath,amssymb,amsthm,mathtools}
\usepackage[colorlinks=true,linkcolor=blue,citecolor=blue]{hyperref}

\newtheorem{theorem}{Theorem}
\newtheorem{proposition}{Proposition}
\newtheorem{lemma}{Lemma}
\theoremstyle{remark}
\newtheorem*{remark}{Remark}

\newcommand{\dd}{\,\mathrm{d}}
\newcommand{\R}{\mathbb{R}}
\DeclareMathOperator{\Var}{Var}

\title{\textbf{Channel~B, tightened: exact asymptotics for the SNIC phase channel}\\[2pt]
\large the $2/3$ exponent proved, the constant identified, the universal object named}
\author{Solomon Williams \\ \small University of Edinburgh}
\date{\today}

\begin{document}
\maketitle

\begin{abstract}\noindent
We raise Channel~B to the rigour of Channel~A. For the noisy saddle-node-on-circle (the
SNIC inner phase equation) the noise-induced rotation and phase diffusion are governed by
the \emph{exact} mean-first-passage-time quadrature of a $1$D diffusion; Watson's lemma on
its degenerate-saddle asymptotics \emph{proves} the $2/3$ exponent and identifies the
constants as definite integrals over the parameter-free canonical noisy saddle-node
$\dd u=\tfrac12 u^2\dd\tau+\dd W$:
\[
  \omega=\frac{\pi}{J}\,\sigma^{2/3}(1+o(1)),\qquad
  D_\phi=\frac{\pi^2 V}{4J^3}\,\sigma^{2/3}(1+o(1)),
\]
with $J=\iint_{w<u}e^{(w^3-u^3)/3}\dd w\dd u$ ($\langle T_{\rm canon}\rangle=2J$) and
$V=\Var(T_{\rm canon})$. This \emph{proves} the exponent $2/3$ for both observables,
resolving the earlier pre-asymptotic ``$0.83$'' (a finite-$\sigma$ diffusion estimate),
and names the universal object---the canonical noisy saddle-node first-passage law---as the
phase-channel analogue of Tracy--Widom. Numerically $J\approx5$, $\pi/J\approx0.6$,
consistent with the measured $\omega/\sigma^{2/3}\approx0.627$.
\end{abstract}

\section{The exact quadrature}

The SNIC inner phase dynamics at criticality is the noisy Adler bottleneck
$\dd\theta=(1-\cos\theta)\dd t+\sigma\dd W$ on $S^1$ (period $2\pi$), a $1$D diffusion with
drift $f(\theta)=1-\cos\theta=-U'(\theta)$, $U(\theta)=\sin\theta-\theta$. Its mean
velocity is exact \cite{Risken}:
\begin{equation}\label{eq:vel}
  \omega=\langle\dot\theta\rangle=\frac{2\pi}{\langle T\rangle},\qquad
  \langle T\rangle=\frac{2}{\sigma^2}\int_{0}^{2\pi}\!\!\dd y\,e^{2U(y)/\sigma^2}
       \int_{y-2\pi}^{y}\!\!\dd z\,e^{-2U(z)/\sigma^2},
\end{equation}
the mean first-passage time to advance one period (Pontryagin/scale-function formula). The
SNIC is the degenerate point $\theta=0$ where $f$ vanishes \emph{quadratically}
($f\approx\tfrac12\theta^2$, $U\approx-\theta^3/6$): a saddle-node in the phase.

\section{The $2/3$ exponent, proved}

\begin{proposition}[exact criticality asymptotics]\label{prop:23}
As $\sigma\to0$,
\[
  \omega=\frac{\pi}{J}\,\sigma^{2/3}\,(1+o(1)),\qquad
  J=\iint_{w<u}e^{(w^3-u^3)/3}\dd w\,\dd u,
\]
and the phase diffusion $D_\phi=\dfrac{\pi^2 V}{4J^3}\,\sigma^{2/3}(1+o(1))$ with
$V=\Var(T_{\rm canon})$. Both exponents are exactly $2/3$.
\end{proposition}

\begin{proof}
The integral \eqref{eq:vel} concentrates at the degenerate saddle $\theta=0$, where
$2U/\sigma^2\approx-\theta^3/(3\sigma^2)$. Substituting $\theta=\sigma^{2/3}u$,
$z=\sigma^{2/3}w$ (so $2U/\sigma^2\mapsto-u^3/3$ and $\dd y\,\dd z=\sigma^{4/3}\dd u\dd w$)
turns the leading part of \eqref{eq:vel} into
\[
  \langle T\rangle=\frac{2}{\sigma^2}\,\sigma^{4/3}\!\iint_{w<u}\!e^{(w^3-u^3)/3}\dd w\,\dd u\,(1+o(1))
                 =2J\,\sigma^{-2/3}(1+o(1)),
\]
a rigorous Watson-lemma estimate (the integrand has integrable $1/u^2$ tails, the
sub-leading non-bottleneck time is $O(1)$). Hence $\omega=2\pi/\langle T\rangle=
(\pi/J)\sigma^{2/3}(1+o(1))$. The same rescaling is the change of variables
$\theta=\sigma^{2/3}u$, $t=\sigma^{-2/3}\tau$ carrying the local bottleneck
$\dd\theta=\tfrac12\theta^2\dd t+\sigma\dd W$ to the \emph{parameter-free canonical noisy
saddle-node} $\dd u=\tfrac12 u^2\dd\tau+\dd W$, whose first-passage time has mean $2J$ and
variance $V$. As one rotation is one bottleneck passage, the rotation is a renewal process
of rate $\omega$ and step $2\pi$; its diffusion is
$D_\phi=(2\pi)^2\,\Var(T)/2\langle T\rangle^3$, and inserting
$\langle T\rangle=2J\sigma^{-2/3}$, $\Var(T)=V\sigma^{-4/3}$ gives
$D_\phi=(\pi^2 V/4J^3)\sigma^{2/3}(1+o(1))$.
\end{proof}

\begin{remark}[the universal object]
$\dd u=\tfrac12 u^2\dd\tau+\dd W$ is to Channel~B what the stochastic Airy ground state is
to Channel~A: a \emph{parameter-free} canonical inner object whose first-passage law is the
universal phase-passage distribution. Channel~A's universal law (Tracy--Widom) is closed
form; Channel~B's is defined by this canonical SDE, with $\langle T\rangle=2J$,
$\Var=V$ definite numbers.
\end{remark}

\section{Resolving the ``$0.83$''}

Proposition~\ref{prop:23} \emph{proves} that the diffusion exponent is $2/3$, identical to
the drift. The previously reported $0.83$ (and a re-measured $0.78$) is the
\emph{finite-$\sigma$ effective} diffusion exponent: near criticality $D_\phi$ carries a
slowly-decaying correction (the constant $V$ has a broad first-passage distribution,
$\mathrm{CV}\approx0.6$), so the log-log slope approaches $2/3$ from above. It was never a
distinct exponent; the rigorous value is $2/3$.

\section{Numerics}

\begin{itemize}
\item $J=\iint_{w<u}e^{(w^3-u^3)/3}\approx 5.0$ (slowly-converging algebraic tails;
$5.24$ at cutoff $15$, $5.41$ at $20$), giving $\pi/J\approx0.60$, against the simulated
$\omega/\sigma^{2/3}\approx0.627$ at $\sigma\sim0.15$.
\item The canonical first-passage time $\langle T_{\rm canon}\rangle\to2J$: simulating
$\dd u=\tfrac12 u^2\dd\tau+\dd W$ from $u=-6$ gives $\langle T\rangle=9.27$ (rising toward
$2J\approx10.8$ as the start recedes), $\Var\approx33$.
\item The drift exponent is measured at $\sigma^{0.668}$ ($\to2/3$); the diffusion's
effective exponent $0.73$--$0.78$ drifts toward $2/3$ as $\sigma\to0$, as predicted.
\end{itemize}

\section{Scope: now level with Channel~A}

\begin{center}
\begin{tabular}{lll}
\hline
 & Channel~A (amplitude) & Channel~B (SNIC phase) \\
\hline
inner object & stochastic Airy / $\dd R=(R^2-Y)\dd T+\eta\dd B$ & noisy saddle-node-on-circle\\
canonical form & $-g''+xg=g^3$ (Painlev\'e~II) & $\dd u=\tfrac12 u^2\dd\tau+\dd W$\\
universal law & Tracy--Widom (closed form) & canonical first-passage law\\
exponent & $\tfrac12$ (proved) & $\tfrac23$ (proved, Prop.~\ref{prop:23})\\
constant & $c=\int g'^2=5.4439\ldots$ & $\pi/J$, $\;J=\iint_{w<u}e^{(w^3-u^3)/3}$\\
\hline
\end{tabular}
\end{center}

\noindent The structure now matches: an exact asymptotic, a \emph{proved} exponent, the
constant identified as a definite integral, and the universal inner object named. The one
honest gap relative to Channel~A is precision: $c$ was Pohozaev-pinned to seven digits,
whereas $J$'s algebraic tails make its high-precision value harder to fix (it is
$\approx5$); and Channel~B's universal law lacks a Tracy--Widom-style closed form. But the
\emph{rigour} of the exponent and the asymptotic form is now equal to Channel~A's.

\begin{thebibliography}{9}
\bibitem{Risken} H.~Risken, \emph{The Fokker--Planck Equation}, 2nd ed., Springer (1989),
\S11.4 (mean velocity / MFPT for tilted periodic potentials).
\bibitem{Reimann} P.~Reimann et al., \emph{Giant acceleration of free diffusion}, PRL
\textbf{87} (2001) 010602; PRE \textbf{65} (2002) 031104.
\bibitem{int} Companion: \texttt{ChannelB\_SNIC\_proof} (the $2/3$ reduction),
\texttt{InstantonPrinciple\_proof} (Channel~A rate), \texttt{PathA\_endgame\_writeup}.
\end{thebibliography}

\end{document}
